\documentclass[11pt]{article}
% =========================================================================
% SHARED STYLE FILE for 11_Master_Projects reports
% Visual style follows 0_Thesis_Submission/24m0553_AK_Gaurav_Mtech_Report.pdf:
% serif body font, blue hyperlinked TOC/links, ruled section headings,
% IITB-style title page, formal academic report structure.
% =========================================================================

\usepackage[a4paper,margin=1in]{geometry}
\usepackage{xcolor}
\Sepackage{graphicx}
graphicspath{{../images/}}
\usepackage{amsmath,amssymb}
\usepackage{booktabs}
\usepackage{caption}
\usepackage{float}
\usepackage[hidelinks]{hyperref}
\usepackage{titlesec}
\usepackage{enumitem}
\usepackage{setspace}
\usepackage[T1]{fontenc}
\usepackage{lmodern}
\usepackage{microtype}

% Every \includegraphics call is automatically framed with a thin black
% border (matches the printed look of figures/photos/charts in the source
% reports). \oldincludegraphics is kept available for institutional
% logos/seals on title pages, which should remain unframed.
\let\oldincludegraphics\includegraphics
\setlength{\fboxsep}{1pt}
\setlength{\fboxrule}{0.6pt}
\renewcommand{\includegraphics}[2][]{\fbox{\oldincludegraphics[#1]{#2}}}

\definecolor{thesisblue}{RGB}{0,0,205}
\hypersetup{
  colorlinks=true,
  linkcolor=thesisblue,
  urlcolor=thesisblue,
  citecolor=thesisblue,
  filecolor=thesisblue
}

\onehalfspacing
\setlength{\parskip}{6pt}
\setlength{\parindent}{0pt}

\titleformat{\section}
  {\normalfont\Large\bfseries}{\thesection}{1em}{}
\titleformat{\subsection}
  {\normalfont\large\bfseries}{\thesubsection}{1em}{}
\titleformat{\subsubsection}
  {\normalfont\normalsize\bfseries}{\thesubsubsection}{1em}{}

\newcommand{\reportTitlePage}[7]{%
  % #1 = Report title
  % #2 = Course code & name (e.g. "CE 605: Applied Statistics")
  % #3 = Subtitle/topic line (e.g. "Project Report 2: Network Analysis")
  % #4 = Author block (can be multi-line with \\ for group projects)
  % #5 = Supervisor name(s)
  % #6 = Programme line (e.g. "M.Tech in Transportation Systems Engineering")
  % #7 = Month/Year
  \begin{titlepage}
    \centering
    \vspace*{0.5cm}
    {\large\bfseries #2}\par
    \vspace{0.3cm}
    {\Large\bfseries #3}\par
    \vspace{1.2cm}
    {\Huge\bfseries #1}\par
    \vspace{1.5cm}
    \oldincludegraphics[width=3.2cm]{../../assets/iitb_logo.png}\par
    \vspace{1.2cm}
    Submitted by:\par
    \vspace{0.2cm}
    {\bfseries #4}\par
    \vspace{0.8cm}
    Under the Supervision of\par
    \vspace{0.2cm}
    {\bfseries #5}\par
    \vspace{1.2cm}
    #6\par
    Department of Civil Engineering\par
    \textbf{Indian Institute of Technology Bombay}\par
    Mumbai -- 400076, India\par
    \vspace{0.3cm}
    #7\par
  \end{titlepage}
}

\newcommand{\reportAbstract}[1]{%
  \section*{Abstract}
  #1
  \vspace{0.5cm}
}


\title{Economic Evaluation of Delhi Metro Phase 4A}
\author{Group 30}
\date{}

\begin{document}

\reportTitlePage
  {Economic Evaluation of Delhi Metro Phase 4A}
  {CE 776: Transportation Project Evaluation and Decision Making}
  {Course Project, Group 30}
  {Lekshmi Muraleedharan (24M0544) \\ Yedukrishna K R (24M0545) \\ Sanjeev Patel (24M0546) \\
   Ankit Raj (24M0548) \\ Rajasri S (24M0549) \\ Abhijeet Kumar Gaurav (24M0553)}
  {Prof. Perumal Vedagiri}
  {M.Tech in Transportation Systems Engineering}
  {Spring 2025}

\pagenumbering{roman}

\reportAbstract{
This report presents an economic evaluation of Delhi Metro Phase~4A, a 61.679~km, three-corridor
mixed underground-and-elevated expansion of the Delhi Metro network (Aerocity--Tughlakabad,
Janakpuri West--R.K.\ Ashram, and Mukundpur--Maujpur), including one of the network's longest
underground corridors. Using standard project-appraisal techniques, namely Cost-Benefit Analysis (CBA),
Economic Internal Rate of Return (EIRR), Net Present Value (NPV), and Benefit-Cost Ratio (BCR), the
project's economic viability is assessed against its estimated completion cost of Rs.~24,948.65~crore.
The evaluation finds an EIRR of 19.81\% (without discounting) and 6.97\% (at a 12\% discount rate),
both exceeding or closely approaching the discount-rate threshold; a positive NPV of Rs.~394,074.7~crore
(undiscounted) and Rs.~26,136.2~crore (discounted); and a Benefit-Cost Ratio of 5.7 (undiscounted) and
1.99 (discounted), indicating that discounted benefits are nearly twice the discounted costs even
under conservative assumptions. Sensitivity analysis confirms the project remains viable across
tested scenarios, with EIRR never falling below 6.97\% and the discounted BCR staying above 1.99. The
evaluation concludes that Delhi Metro Phase~4A is economically justified, offering a cost-effective and
sustainable response to Delhi's growing urban transit demand.

\vspace{0.3cm}
\noindent\textbf{Keywords:} Cost-Benefit Analysis; Economic Internal Rate of Return; Net Present Value;
Benefit-Cost Ratio; Metro Rail Appraisal; Delhi Metro Phase 4
}

\newpage
\tableofcontents
\newpage
\pagenumbering{arabic}

\section{Project Overview}

\subsection{Project Description}
Delhi Metro Phase~4A is a mixed underground-and-elevated metro expansion comprising one of the longest
underground corridors in the Delhi Metro network. Table~\ref{tab:corridors} summarises the three
corridors that make up the project.

\begin{table}[H]
  \centering
  \caption{Corridor-wise route length for Delhi Metro Phase 4A.}
  \label{tab:corridors}
  \begin{tabular}{@{}lrrr@{}}
    \toprule
    \textbf{Corridor} & \textbf{Total (km)} & \textbf{Underground (km)} & \textbf{Elevated (km)} \\
    \midrule
    Aerocity -- Tughlakabad       & 20.201 & 14.619 & 5.582 \\
    Janakpuri West -- R.K.\ Ashram & 28.920 & 7.740  & 21.180 \\
    Mukundpur -- Maujpur          & 12.558 & 0.000  & 12.558 \\
    \midrule
    \textbf{Total}                & \textbf{61.679} & \textbf{22.359} & \textbf{39.320} \\
    \bottomrule
  \end{tabular}
\end{table}

Across the three corridors, the project adds 79 stations (17 underground, 62 at/above grade,
Table~\ref{tab:stations}) and is designed for a maximum speed of 95~km/h with a scheduled operating
speed of 35~km/h.

\begin{table}[H]
  \centering
  \caption{Corridor-wise station count.}
  \label{tab:stations}
  \small
  \begin{tabular}{@{}p{4.3cm}rrrr@{}}
    \toprule
    \textbf{Corridor} & \textbf{Length (km)} & \textbf{UG stations} & \textbf{Elevated stations} & \textbf{Total} \\
    \midrule
    Aerocity -- Saket -- Tughlakabad & 20.20 & 10 & 5  & 15 \\
    Janakpuri (West) -- R.K.\ Ashram & 28.92 & 7  & 18 & 25 \\
    Mukundpur -- Burari -- Maujpur   & 12.54 & 0  & 6  & 6 \\
    \midrule
    \textbf{Total}                    & \textbf{61.66} & \textbf{17} & \textbf{29} & \textbf{79} \\
    \bottomrule
  \end{tabular}
\end{table}

\subsection{Purpose}
The project is intended to (i) ease traffic pressure on major roads such as the Mehrauli-Badarpur Road
and NH-48 and enhance connectivity between key urban locations and the airport; (ii) serve high
ridership potential, with projected daily demand reaching 19.4~lakh passengers by 2026 (1.8~million by
2031 and 2.59~million by 2041); and (iii) support sustainable urban transport by shifting travel to
mass transit, thereby reducing carbon emissions, promoting energy-efficient travel, and improving air
quality in Delhi.

\begin{figure}[H]
  \centering
  \includegraphics[width=0.55\textwidth]{fig_metro_train.png}
  \caption{Delhi Metro Rail Corporation rolling stock.}
\end{figure}

\subsection{Construction Methodology}
Elevated sections use viaducts of prestressed-concrete ``box'' girders on single piers with pile or
open foundations, while the underground section is built using tunnel boring machines, with cut-and-cover
construction at underground stations.

\subsection{Current Status and Funding}
At the time of this evaluation, 2.03~km of Phase~4A was operational, 63.07~km was under construction,
47.22~km was approved, and 10.17~km remained on hold/proposed. Funding sources for the project include
the Government of India, the Delhi Development Authority (DDA), the Delhi Government, private
investors, JICA, and domestic financial institutions.

\section{Project Cost}
The total estimated cost of the three corridors plus additional rolling stock, at January-2019 price
levels including all taxes, is Rs.~22,493~crore; including Interest During Construction (IDC), the
completion cost rises to Rs.~24,948.65~crore (Table~\ref{tab:cost}).

\begin{table}[H]
  \centering
  \caption{Corridor-wise estimated completion cost (Rs.\ crore).}
  \label{tab:cost}
  \small
  \begin{tabular}{@{}p{4.5cm}p{3.7cm}p{3.2cm}@{}}
    \toprule
    \textbf{Corridor} & \textbf{Estimated cost (Jan-2019 level, incl.\ taxes)} & \textbf{Completion cost (excl.\ IDC)} \\
    \midrule
    Aerocity -- Tughlakabad        & 8,231  & 9,217 \\
    Janakpuri West -- R.K.\ Ashram & 9,349  & 10,357 \\
    Mukundpur -- Maujpur           & 2,190  & 2,456 \\
    Additional rolling stock       & 2,723  & 2,859 \\
    \midrule
    \textbf{Total}                 & \textbf{22,493} & \textbf{24,889} \\
    \bottomrule
  \end{tabular}
\end{table}

\begin{table}[H]
  \centering
  \caption{Year-wise investment schedule (Rs.\ crore).}
  \small
  \begin{tabular}{@{}p{2.0cm}p{2.6cm}p{2.6cm}p{1.3cm}p{2.6cm}@{}}
    \toprule
    \textbf{Year} & \textbf{Est.\ constr.\ cost} & \textbf{Completion cost (taxes)} & \textbf{IDC} & \textbf{Completion cost (taxes+IDC)} \\
    \midrule
    2019--20 & 987.99  & 987.99  & --    & 987.99 \\
    2020--21 & 5688.53 & 5954.61 & 4.66  & 5959.27 \\
    2021--22 & 4349.36 & 4730.46 & 9.06  & 4739.52 \\
    2022--23 & 5930.95 & 6710.46 & 18.87 & 6729.33 \\
    2023--24 & 5535.55 & 6505.62 & 26.92 & 6532.54 \\
    \midrule
    \textbf{Total} & \textbf{22492.38} & \textbf{24889.14} & \textbf{59.51} & \textbf{24,948.65} \\
    \bottomrule
  \end{tabular}
\end{table}

\section{Economic Evaluation Framework}

\subsection{Cost-Benefit Analysis}
The economic evaluation follows standard Cost-Benefit Analysis (CBA) methodology, comparing the
project's economic costs against its economic benefits over its appraisal life.

\textbf{Costs included:} capital (construction) cost, land acquisition cost, and operating and
maintenance costs, along with rolling-stock, environmental, utility-relocation, and contingency costs
and interest during construction.

\textbf{Benefits included:} travel-time savings, vehicle operating cost (VOC) savings, reduced air
pollution and emissions, reduction in road accidents, and fuel savings. These are summarised as annual
savings in road-infrastructure maintenance costs, annual fuel savings for metro passengers and diverted road
users, annual travel-time savings for metro passengers and (via decongestion) for remaining road users,
annual reduction in vehicle operating costs, cost savings from reduced emissions, and reduction in
accident-related costs.

\subsection{Evaluation Indicators}
Three principal indicators were used to assess project viability:
\begin{itemize}
  \item \textbf{Economic Internal Rate of Return (EIRR):} the discount rate at which the present
  value of economic benefits equals the present value of economic costs; if EIRR exceeds the assumed
  discount rate, the project is economically viable.
  \item \textbf{Net Present Value (NPV):} the difference between the present value of benefits and
  the present value of costs over the project's life; a positive NPV indicates the project adds
  economic value to society.
  \item \textbf{Benefit-Cost Ratio (BCR):} the ratio of discounted benefits to discounted costs; a
  BCR greater than 1 indicates benefits exceed costs.
\end{itemize}

A sensitivity analysis was additionally performed to assess how changes in traffic and cost assumptions
affect project viability, using EIRR and BCR as the response indicators.

\section{Economic Evaluation Results}

\subsection{Base-Case Indicators}
Table~\ref{tab:econindicators} reports the base-case (2048--49 appraisal-period) economic indicator
values, with and without discounting.

\begin{table}[H]
  \centering
  \caption{Economic indicator values (base case, 2048--49).}
  \label{tab:econindicators}
  \begin{tabular}{@{}lrr@{}}
    \toprule
    & \textbf{Without discount} & \textbf{With discount (12\%)} \\
    \midrule
    Total cumulative cost (Rs.\ crore)    & 83,707.27  & 26,397.66 \\
    Total cumulative benefit (Rs.\ crore) & 477,782    & 52,534 \\
    Benefit-Cost Ratio                    & 5.70       & 1.99 \\
    Net Present Value (Rs.\ crore)        & 394,074.7  & 26,136.2 \\
    EIRR                                  & 19.81\%    & 6.97\%\textsuperscript{*} \\
    \bottomrule
  \end{tabular}
  \\[2pt]
  \raggedright\footnotesize\textsuperscript{*}A companion slide in the source deck reports 6.91\% for
  the same discounted-EIRR figure; the discrepancy is a minor rounding/transcription difference in the
  original material and does not affect the conclusion that EIRR remains well above typical hurdle
  rates for public transit investment.
\end{table}

Since EIRR (19.81\% undiscounted) exceeds the 12\% discount rate used, and the NPV is positive under
both undiscounted and discounted assumptions, the project is economically viable: for every Rs.~1
invested, the undiscounted analysis projects Rs.~5.7 of benefits, and even under the more conservative
12\%-discounted analysis, benefits are projected at nearly double the costs (BCR~=~1.99).

A separate Financial Internal Rate of Return (FIRR) was also estimated: 8.03\% without Transit-Oriented
Development (TOD) and Value Capture Financing (VCF) income, rising to 9.60\% when TOD/VCF revenue
streams are included, indicating that non-fare revenue mechanisms materially improve the project's
financial (as distinct from economic) return.

\subsection{Sensitivity Analysis}
Sensitivity analysis assessed how changes in traffic and cost assumptions affect project viability. The
project remained viable across the scenarios tested, with EIRR never falling below 6.97\% and the
discounted Benefit-Cost Ratio never falling below 1.99, indicating that the economic case for Phase~4A
is robust to reasonable variation in underlying assumptions.

\section{Project Benefits}
Project benefits are classified as direct and indirect (Table~\ref{tab:benefits}).

\begin{table}[H]
  \centering
  \caption{Direct and indirect project benefits.}
  \label{tab:benefits}
  \begin{tabular}{@{}ll@{}}
    \toprule
    \textbf{Direct benefits} & \textbf{Indirect benefits} \\
    \midrule
    Expedited commutes       & Job creation \\
    Reduced congestion       & Business growth \\
    Vehicle cost savings     & Greener environment \\
    Economical transport     & Property appreciation \\
    Enhanced safety          & Improved air quality \\
    Fare revenue             & Economic development \\
    \bottomrule
  \end{tabular}
\end{table}

\begin{figure}[H]
  \centering
  \includegraphics[width=0.85\textwidth]{fig_metro_network_map.png}
  \caption{Delhi Metro network, showing Phase I--IV and proposed monorail corridors.}
\end{figure}

\section{Conclusion and Recommendations}

\subsection{Economic Viability Summary}
\begin{itemize}
  \item \textbf{EIRR (19.81\%)} exceeds the 12\% discount-rate threshold, indicating strong
  socio-economic benefits such as reduced travel time, fewer accidents, and environmental gains.
  \item \textbf{NPV is positive} under both discounted and undiscounted assumptions, reinforcing the
  project's long-term economic viability.
  \item \textbf{Benefit-Cost Ratio} of 5.71 (undiscounted) and 1.99 (at 12\% discount) shows that even
  under conservative assumptions, benefits substantially exceed costs.
  \item \textbf{Sensitivity analysis} confirms the project remains viable across scenarios tested,
  with EIRR never falling below 6.97\% and discounted BCR staying above 1.99.
\end{itemize}

\subsection{Key Takeaways for Decision-Makers}
A strong EIRR and BCR confirm the economic viability of the project; earlier completion yields higher
cumulative benefits (as delay defers benefit realisation while cost escalation continues); the project
is critical for sustainable urban growth in the National Capital Region; the corridor alignment
strategically links major employment, residential, and transit nodes; and there is an opportunity for
policy synergies through last-mile connectivity investment alongside the metro expansion.

\subsection{Learnings}
This exercise provided hands-on experience in economic evaluation techniques, namely Cost-Benefit
Analysis, EIRR, and NPV, and in interpreting sensitivity analysis to assess project risk under
uncertainty. It
also highlighted the real-world impact of large transit infrastructure investment on urban mobility
and the importance of data-driven planning in public transport project appraisal.

\subsection{Concluding Remarks}
Delhi Metro Phase~4A is economically justified, not only meeting but significantly surpassing standard
evaluation thresholds. It stands as a cost-effective, sustainable response to Delhi's growing urban
transit demand, illustrating how thoughtful transport planning can deliver economic, environmental, and
social benefits simultaneously, while underscoring the need for inter-agency coordination and timely
execution to realise its full potential.

\begin{thebibliography}{9}
\bibitem{dmrc} Delhi Metro Rail Corporation Ltd. \textit{Detailed Project Report: Delhi Metro Phase 4}.
\bibitem{mort} Ministry of Housing and Urban Affairs, Government of India. \textit{Metro Rail Policy}.
\end{thebibliography}

\end{document}
